%beamer

% Start with video if you like it and can get the sound working well:
% https://www.youtube.com/watch?v=Dy0hJWltsyE

% Comment/uncomment this line to toggle handout mode
%\newcommand{\handout}{}

% % \bigtimes abgeschrieben von http://tex.stackexchange.com/questions/14386/importing-a-single-symbol-from-a-different-font
% \DeclareFontFamily{U}{mathx}{\hyphenchar\font45}
% \DeclareFontShape{U}{mathx}{m}{n}{
%       <5> <6> <7> <8> <9> <10> gen * mathx
%       <10.95> mathx10 <12> <14.4> <17.28> <20.74> <24.88> mathx12
%       }{}
% \DeclareSymbolFont{mathx}{U}{mathx}{m}{n}
% \DeclareMathSymbol{\bigtimes}{\mathop}{mathx}{161}

\RequirePackage{xcolor}

\def\9{\square}
%\def\9{\blank}

% f"ur Aussagenlogik
\colorlet{alcolor}{blue}
\RequirePackage{tikz}
\usetikzlibrary{arrows.meta}
\newcommand{\alimpl}{\mathrel{\tikz[x={(0.1ex,0ex)},y={(0ex,0.1ex)},>={Classical TikZ Rightarrow[]}]{\draw[alcolor,->,line width=0.7pt,line cap=round] (0,0) -- (15,0);\path (0,-6);}}}
\newcommand{\aleqv}{\mathrel{\tikz[x={(0.1ex,0ex)},y={(0ex,0.1ex)},>={Classical TikZ Rightarrow[]}]{\draw[alcolor,<->,line width=0.7pt,line cap=round] (0,0) -- (18,0);\path (0,-6);}}}
\newcommand{\aland}{\mathbin{\raisebox{-0.6pt}{\rotatebox{90}{\texttt{\color{alcolor}\char62}}}}}
\newcommand{\alor}{\mathbin{\raisebox{-0.8pt}{\rotatebox{90}{\texttt{\color{alcolor}\char60}}}}}
%\newcommand{\ali}[1]{_{\mathtt{\color{alcolor}#1}}}
\newcommand{\alv}[1]{\mathtt{\color{alcolor}#1}}
\newcommand{\alnot}{\mathop{\tikz[x={(0.1ex,0ex)},y={(0ex,0.1ex)}]{\draw[alcolor,line width=0.7pt,line cap=round,line join=round] (0,0) -- (10,0) -- (10,-4);\path (0,-8) ;}}}
\newcommand{\alP}{\alv{P}} %ali{#1}}
%\newcommand{\alka}{\negthinspace\hbox{\texttt{\color{alcolor}(}}}
\newcommand{\alka}{\negthinspace\text{\texttt{\color{alcolor}(}}}
%\newcommand{\alkz}{\texttt{\color{alcolor})}}\negthinspace}
\newcommand{\alkz}{\text{\texttt{\color{alcolor})}}\negthinspace}
\newcommand{\AAL}{A_{AL}}
\newcommand{\LAL}{\hbox{\textit{For}}_{AL}}
\newcommand{\AxAL}{\hbox{\textit{Ax}}_{AL}}
\newcommand{\AxEq}{\hbox{\textit{Ax}}_{Eq}}
\newcommand{\AxPL}{\hbox{\textit{Ax}}_{PL}}
\newcommand{\AALV}{\hbox{\textit{Var}}_{AL}}
\newcommand{\MP}{\hbox{\textit{MP}}}
\newcommand{\GEN}{\hbox{\textit{GEN}}}
\newcommand{\W}{\ensuremath{\hbox{\textbf{w}}}\xspace}
\newcommand{\F}{\ensuremath{\hbox{\textbf{f}}}\xspace}
\newcommand{\WF}{\ensuremath{\{\W,\F\}}\xspace}
\newcommand{\val}{\hbox{\textit{val}}}
\newcommand{\valDIb}{\val_{D,I,\beta}}

\newcommand*{\from}{\colon}

% die nachfolgenden Sachen angepasst an cmtt
\newlength{\ttquantwd}
\setlength{\ttquantwd}{1ex}
\newlength{\ttquantht}
\setlength{\ttquantht}{6.75pt}
\def\plall{%
  \tikz[line width=0.67pt,line cap=round,line join=round,baseline=(B),alcolor] {
    \draw (-0.5\ttquantwd,\ttquantht) -- node[coordinate,pos=0.4] (lll){} (-0.25pt,-0.0pt) -- (0.25pt,-0.0pt) -- node[coordinate,pos=0.6] (rrr){} (0.5\ttquantwd,\ttquantht);
    \draw (lll) -- (rrr);
    \coordinate (B) at (0,-0.35pt);
  }%
}
\def\plexist{%
  \tikz[line width=0.67pt,line cap=round,line join=round,baseline=(B),alcolor] {
    \draw (-0.9\ttquantwd,\ttquantht) -- (0,\ttquantht) -- node[coordinate,pos=0.5] (mmm){} (0,0) --  (-0.9\ttquantwd,0);
    \draw (mmm) -- ++(-0.75\ttquantwd,0);
    \coordinate (B) at (0,-0.35pt);
  }\ensuremath{\,}%
}
\let\plexists=\plexist
\newcommand{\NT}[1]{\ensuremath{\langle\mathrm{#1} \rangle}}

\newcommand{\CPL}{\text{\itshape Const}_{PL}}
\newcommand{\FPL}{\text{\itshape Fun}_{PL}}
\newcommand{\RPL}{\text{\itshape Rel}_{PL}}
\newcommand{\VPL}{\text{\itshape Var}_{PL}}
\newcommand{\ATer}{A_{\text{\itshape Ter}}}
\newcommand{\ARel}{A_{\text{\itshape Rel}}}
\newcommand{\AFor}{A_{\text{\itshape For}}}
\newcommand{\LTer}{L_{\text{\itshape Ter}}}
\newcommand{\LRel}{L_{\text{\itshape Rel}}}
\newcommand{\LFor}{L_{\text{\itshape For}}}
\newcommand{\NTer}{N_{\text{\itshape Ter}}}
\newcommand{\NRel}{N_{\text{\itshape Rel}}}
\newcommand{\NFor}{N_{\text{\itshape For}}}
\newcommand{\PTer}{P_{\text{\itshape Ter}}}
\newcommand{\PRel}{P_{\text{\itshape Rel}}}
\newcommand{\PFor}{P_{\text{\itshape For}}}

\newcommand{\plka}{\alka}
\newcommand{\plkz}{\alkz}
%\newcommand{\plka}{\plfoo{(}}
%\newcommand{\plkz}{\plfoo{)}}
\newcommand{\plcomma}{\hbox{\texttt{\color{alcolor},}}}
\newcommand{\pleq}{{\color{alcolor}\,\dot=\,}}

% MODIFIED (DJ)
% previously: \newcommand{\plfoo}[1]{\mathtt{\color{alcolor}#1}}
\newcommand{\plfoo}[1]{\texttt{\color{alcolor}#1}}

\newcommand{\plc}{\plfoo{c}}
\newcommand{\pld}{\plfoo{d}}
\newcommand{\plf}{\plfoo{f}}
\newcommand{\plg}{\plfoo{g}}
\newcommand{\plh}{\plfoo{h}}
\newcommand{\plx}{\plfoo{x}}
\newcommand{\ply}{\plfoo{y}}
\newcommand{\plz}{\plfoo{z}}
\newcommand{\plR}{\plfoo{R}}
\newcommand{\plS}{\plfoo{S}}

\newcommand{\bv}{\mathrm{bv}}
\newcommand{\fv}{\mathrm{fv}}

%\newcommand{\AxAL}{\hbox{\textit{Ax}}_{AL}}
%\newcommand{\AALV}{\hbox{\textit{Var}}_{AL}}

%\renewcommand{\#}[1]{\literal{#1}}
\newcommand{\A}{\mathcal{A}}
\newcommand{\Adr}{\text{Adr}}
\newcommand{\ar}{\mathrm{ar}}
\newcommand{\ascii}[1]{\literal{\char#1}}
%\newcommand{\assert}[1]{\text{/\!\!/\ } #1}
\newcommand{\assert}[1]{\colorbox{black!7!white}{\ensuremath{\{\;#1\;\}}}}
\newcommand{\Assert}[1]{$\langle$\textit{#1}$\rangle$}
\newcommand{\B}{\mathcal{B}}
\newcommand{\bfmod}{\mathbin{\kw{ mod }}}
\newcommand{\bb}{{\text{bb}}}
\def\bottom{\hbox{\small$\pmb{\bot}$}}
\newcommand{\card}[1]{|#1|}
%\newcommand{\cod}{\mathop{\text{cod}}}  % ist in thwmathabbrevs
\newcommand{\Conf}{\mathcal{C}}
\newcommand{\define}[1]{\emph{#1}}
%\renewcommand{\dh}{d.\,h.\@\xspace}
%\newcommand{\Dh}{D.\,h.\@\xspace}
%\newcommand{\engl}[1]{engl.\xspace\emph{#1}}
\newcommand{\eps}{\varepsilon}
%\newcommand{\evtl}{evtl.\@\xspace}
\newcommand{\fbin}{\text{bin}}
\newcommand{\finv}{\text{inv}}
\newcommand{\fnum}{\text{num}}
\newcommand{\fNum}{{\text{Num}}}
\newcommand{\frepr}{\text{repr}}
\newcommand{\fRepr}{\text{Repr}}
\newcommand{\fZkpl}{\text{Zkpl}}
\newcommand{\fLen}{\text{Len}}
\newcommand{\fsem}{\text{sem}}
\providecommand{\fspace}{\mathord{\text{space}}}
\providecommand{\fSpace}{\mathord{\text{Space}}}
\providecommand{\ftime}{\mathord{\text{time}}}
\providecommand{\fTime}{\mathord{\text{Time}}}
\newcommand{\fTrans}{\text{Trans}}
\newcommand{\fVal}{\text{Val}}

% MODIFIED (DJ)
\newcommand{\Val}{\text{Val}}

%\def\G{\mathbb{Z}}
\newcommand{\HT}[1]{\normalfont\textsc{HT-#1}}
\newcommand{\htr}[3]{\{#1\}\;#2\; \{#3\}}
\newcommand{\Id}{\text{I}}
%\newcommand{\ie}{i.\,e.\@\xspace}
\newcommand{\instr}[2]{\texttt{#1}\ \textit{#2}}
\newcommand{\Instr}[2]{\texttt{#1}\ \textrm{#2}}
\newcommand{\instrr}[3]{\texttt{#1}\ \textit{#2}\texttt{(#3)}}
\newcommand{\Instrr}[3]{\texttt{#1}\ \textrm{#2}\texttt{(#3)}}
\newcommand{\io}{\!\mid\!}
\usepackage{KITcolors}
\newcommand{\literal}[1]{\hbox{\textcolor{blue!95!white}{\textup{\texttt{\scalebox{1.11}{#1}}}}}}
%\newcommand{\literal}[1]{\hbox{\textcolor{KITblue!80!black}{\textup{\texttt{#1}}}}}
\def\kasten#1{\leavevmode\literal{\setlength{\fboxsep}{1pt}\fbox{\vrule  width 0pt height 1.5ex depth 0.5ex #1}}}
\newcommand{\kw}[1]{\ensuremath{\mathbf{#1}}}
\newcommand{\lang}[1]{\ensuremath{\langle#1\rangle}}
%\newcommand{\maw}{m.\,a.\,w.\@\xspace}
%\newcommand{\MaW}{M.\,a.\,w.\@\xspace}
\newcommand{\mdefine}[2][FOOBAR]{\define{#2}\def\foobar{FOOBAR}\def\optarg{#1}\ifx\foobar\optarg\def\optarg{#2}\fi\graffito{\optarg}}
\newcommand{\meins}{\rotatebox[origin=c]{180}{1}}
\newcommand{\Mem}{\text{Mem}}
\newcommand{\memread}{\text{memread}}
\newcommand{\memwrite}{\text{memwrite}}
\providecommand{\meta}[1]{\ensuremath{\langle}\textit{#1}\ensuremath{\rangle}}
%\newcommand{\N}{\mathbb{N}}
\newcommand{\NP}{\mathbf{NP}}
\newcommand{\Nadd}{N_{\text{add}}}
\newcommand{\Nmult}{N_{\text{mult}}}
% MODIFIED (DJ): added \!, mathcal{O}
\newcommand{\Oh}[1]{\mathcal{O}\!\left(#1\right)}
\newcommand{\Om}[1]{\Omega\!\left(#1\right)}
\newcommand{\personname}[1]{\textsc{#1}}
\newcommand{\regname}[1]{\texttt{#1}}
\newcommand{\mima}{\textsc{Mima}\xspace}
\newcommand{\mimax}{\textsc{Mima-X}\xspace}

\def\Pclass{\text{\bfseries P}}
\def\PSPACE{\text{\bfseries PSPACE}}

\newcommand{\SPush}{\text{push}}
\newcommand{\SPop}{\text{pop}}
\newcommand{\SPeek}{\text{peek}}
\newcommand{\STop}{\text{top}}
\newcommand{\STos}{\text{\itshape tos}}
\newcommand{\SBos}{\text{\itshape bos}}

%\newcommand{\R}{\mathbb{R}}
\newcommand{\Rnullplus}{\R_0^{+}}
\newcommand{\Rplus}{\R_{+}}
\newcommand{\resp}{resp.\@\xspace}
\newcommand{\Sem}{\text{Sem}}
\newcommand{\sgn}{\mathop{\text{sgn}}}
\newcommand{\sqbox}{\mathop{\raisebox{-6.2pt}{\hbox{\hbox to 0pt{$^{^{\sqcap}}$\hss}$^{^{\sqcup}}$}}}}
\newcommand{\sqleq}{\sqsubseteq}
\newcommand{\sqgeq}{\sqsupseteq}
% MODIFIED (DJ): added \!
\newcommand{\Th}[1]{\Theta\!\left(#1\right)}
%\newcommand{\usw}{usw.\@\xspace}
\newcommand{\V}[1]{\hbox{\textit{#1}}}
\newcommand{\x}{\times}
\newcommand{\ZK}{\mathbb{K}}
%\newcommand{\Z}{\mathbb{Z}}
\newcommand{\zB}{z.\,B.\@\xspace}
\newcommand{\ZB}{Z.\,B.\@\xspace}
% \newcommand{\bb}{{\text{bb}}}
% \def\##1{\hbox{\textcolor{darkblue}{\texttt{#1}}}}
% \def\A{\mathcal{A}}
% \newcommand{\0}{\#0}
% \newcommand{\1}{\#1}
% \newcommand{\Obj}{\text{Obj}}
% \newcommand{\start}{\mathop{\text{start}}}
% \newcommand{\compactlist}{\addtolength{\itemsep}{-\parskip}}
% \newcommand{\fval}{\text{val}}
% \newcommand{\lang}[1]{\ensuremath{\langle#1\rangle}}
% \newcommand{\io}{\!\mid\!}
% \def\sqbox{\mathop{\raisebox{-6.2pt}{\hbox{\hbox to 0pt{$^{^{\sqcap}}$\hss}$^{^{\sqcup}}$}}}}
% \def\sqleq{\sqsubseteq}
% \def\sqgeq{\sqsupseteq}
\def\Td{T_{\overline{d}}}
% \newcommand{\csym}[1]{\ensuremath{\#{c}_{\#{\hbox{\scriptsize #1}}}}}
% \newcommand{\F}{\ensuremath{\mathcal{F}}}
% \newcommand{\fsym}[2]{\ensuremath{\#{f}^{\#{\hbox{\scriptsize #1}}}_{\#{\hbox{\scriptsize #2}}}}}
% \newcommand{\rsym}[2]{\ensuremath{\#{R}^{\#{\hbox{\scriptsize #1}}}_{\#{\hbox{\scriptsize #2}}}}}
% \newcommand{\xsym}[1]{\ensuremath{\#{x}_{\#{\hbox{\scriptsize #1}}}}}
% \newcommand{\I}{\mathcal{I}}
% ********************************************************************

\usepackage{../TutTexbib/thwregex}

\newcommand{\nM}{\mathbb{M}}
\newcommand{\nR}{\mathbb{R}}
\newcommand{\nN}{\mathbb{N}}
\newcommand{\nZ}{\mathbb{Z}}
\newcommand{\nQ}{\mathbb{Q}}
\newcommand{\nB}{\mathbb{B}}
\newcommand{\nC}{\mathbb{C}}
\newcommand{\nK}{\mathbb{K}}
\newcommand{\nF}{\mathbb{F}}
\newcommand{\nG}{\mathbb{G}}
\newcommand{\nullel}{\mathcal{O}}
\newcommand{\einsel}{\mathds{1}}
\newcommand{\nP}{\mathbb{P}}
\newcommand{\Pot}{\mathcal{P}}
\renewcommand{\O}{\text{O}}

%\newcommand{\literal}[1]{\textcolor{blue}{\texttt{#1}}}
\newcommand{\realTilde}{\textasciitilde \ }
\newcommand{\setsize}[1]{\; \mid #1 \mid \; }

\definecolor{myRed}{RGB}{255,75,20}
\colorlet{myGreen}{KITpalegreen}

\newcounter{tfqtempcount}
\newcommand{\truefalseQ}[4]{
	\setcounter{tfqtempcount}{#1}
	\addtocounter{tfqtempcount}{1}
	\truefalseQuestion{#1}{\value{tfqtempcount}}{#2}{#3}{#4}
}
\newcommand{\truefalseQuestion}[5]{\item<#1-|handout:#1-> \color<#2-|handout:#2->{#3} #4 \qquad \visible<#2-|handout:#2->{#5}}

\newcommand{\boder}{\ensuremath{\textcolor{blue}{\vee}}\text{} }
\newcommand{\bund}{\ensuremath{\textcolor{blue}{\wedge}}\text{}}
\newcommand{\bimp}{\ensuremath{\textcolor{blue}{\to}}\text{}}
\newcommand{\bnot}{\ensuremath{\textcolor{blue}{\neg}}\text{}}
\newcommand{\bone}{\ensuremath{\textcolor{blue}{1}}\text{}}
\newcommand{\bzero}{\ensuremath{\textcolor{blue}{0}}\text{}}
\newcommand{\bleftBr}{\ensuremath{\textcolor{blue}{(}}\text{}}
\newcommand{\brightBr}{\ensuremath{\textcolor{blue}{)}}\text{}}

\newcommand{\summe}[2]{\sum\limits_{#1}^{#2}}
\newcommand{\limes}[1]{\lim\limits_{#1}}

\newcommand{\numpp}{\advance \value{weeknum} by -2 \theweeknum \advance \value{weeknum} by 2}
\newcommand{\nump}{\advance \value{weeknum} by -1 \theweeknum \advance \value{weeknum} by 1}

\newcommand{\mycomment}[1]{}

\newcounter{abc}
\newenvironment{alist}{
  \begin{list}{(\alph{abc})}{
      \usecounter{abc}\setlength{\leftmargin}{8mm}\setlength{\labelsep}{2mm}
    }
}{\end{list}}

\renewcommand{\mod}{\text{ \bf{mod} }}
\renewcommand{\div}{\text{ \bf{div} }}

%\newcommand{\sgn}{\text{sgn}}

% Toggle Handout mode by including this line before including style_tut:
% \newcommand{\handout}{}

%% Beamer-Klasse im korrekten Modus
\ifdefined \handout
\documentclass[handout]{beamer} % Handout mode
\else
\documentclass{beamer}
\fi


% Das ist der KIT-Stil
%\usepackage{../TutTexbib/beamerthemekit}
\usepackage[deutsch,titlepage0]{../TutTexbib/KIT/beamerthemeKITmod}
\TitleImage[width=\titleimagewd]{../figures/titlepage.jpg}
%\usetheme[deutsch,titlepage0]{KIT}

% Include PDFs
\usepackage{pdfpages}

% Libertine font (Original GBI font)
\usepackage{libertine}
%\renewcommand*\familydefault{\sfdefault}  %% Only if the base font of the document is to be sans serif

% Nicer math symbols
\usepackage{eulervm}
%\usepackage{mathpazo}

%% Deutsche Silbentrennung und Beschriftungen
\usepackage[ngerman]{babel}

%% UTF-8-Encoding
\usepackage[utf8]{inputenc}

\usepackage{csquotes}

%% Bibliotheken für viele mathematische Symbole
\usepackage{amsmath, amsfonts, amssymb}

%% Anzeigetiefe für Inhaltsverzeichnis: 1 Stufe
\setcounter{tocdepth}{1}

%% Schönere Schriften
\usepackage[TS1,T1]{fontenc}

%% Tabellen
\usepackage{array}
\usepackage{multicol}

%% Bibliothek für Graphiken
\usepackage{graphicx}

%% der wird sowieso in jeder Datei gesetzt
%%\graphicspath{{../figures/}}

%% Hyperlinks
\usepackage{hyperref}
% I don't know why, but this works and only includes sections and NOT subsections in the pdf-bookmarks.
\hypersetup{bookmarksdepth=subsection} 

%\usepackage{lmodern}
\usepackage{colortbl}
\usepackage[absolute,overlay]{textpos}
\usepackage{listings}
\usepackage{forloop}
%\usepackage{algorithmic} % PseudoCode package 

\usepackage{tikz}
\usetikzlibrary{matrix}
\usetikzlibrary{arrows.meta}
\usetikzlibrary{automata}
\usetikzlibrary{tikzmark}

% Needed for gbi-macros
\usepackage{xspace}

%%%%%%%%%%%% INHALT %%%%%%%%%%%%%%%%

%% Wochennummer
\newcounter{weeknum}

%% Titelinformationen
\title[GBI Tutorium, Woche \theweeknum]{Grundbegriffe der Informatik \\ Tutorium 36}
\subtitle{Termin \theweeknum \ | \mydate \\ Thassilo Helmold}
\author[Thassilo Helmold]{Thassilo Helmold}
\institute{KIT -- Karlsruher Institut für Technologie}
\date{\mydate}

%% Titel einfügen
\newcommand{\titleframe}{\frame{\titlepage}}

%% Alles starten mit \starttut{X}
\newcommand{\starttut}[1]{\setcounter{weeknum}{#1}\titleframe\frame{\frametitle{Inhalt}\tableofcontents} \AtBeginSection[]{%
\begin{frame}
	\tableofcontents[currentsection]
\end{frame}\addtocounter{framenumber}{-1}}}


\newcommand{\framePrevEpisode}{
	\begin{frame}
		\centering
		\textbf{In the previous episode of GBI...}
	\end{frame}
}

%% Legacy: Lastframe. Not for further usage!
\newcommand{\lastframe}[4]{\xkcdframe{#1}{#2}{#3}{#4}{}}
\newcommand{\xkcdframe}[5]{
	\frame[plain]{
		\vspace{-#2pt}
		\begin{figure}[H]
			\centering
			\includegraphics[scale=#1]{#3}
			\vspace{-7pt}
			\caption{ \texttt{\url{#4}} }
			\vspace{5pt}
			#5
		\end{figure} 
	}
}

%% Wörter
\newcommand{\code}[1]{$\mathbf{#1}$}

%% Sterne

\newcounter{starsc}
\newcommand{\stars}[1]{
	\hfill
	\begin{minipage}{100px}
		\forloop{starsc}{0}{\value{starsc} < #1}%
		{%
			\includegraphics[scale=0.05]{star-full.pdf} \hspace*{1px}
		}%
		\forloop{starsc}{\value{starsc}}{\value{starsc} < 5}%
		{%
			\includegraphics[scale=0.05]{star-empty.pdf} \hspace*{1px}
		}
		\vspace*{2px}
	\end{minipage}
}

% No stars for me...
\renewcommand{\stars}[1]{}

\newcommand{\slideThanks}{
	\begin{frame}
		\frametitle{Danksagung}
		\begin{block}{}
			Dieser Foliensatz basiert in Teilen auf Folien von:\\[1em]
			Philipp Basler \\
			Nils Braun \\
			Dominik Doerner \\
			Ou Yue \\
		\end{block}
	\end{frame}
}

%% Verbatim
\usepackage{moreverb}
\graphicspath{{../figures/}}

\begin{document}
\starttut{1}

\section{Organisatorisches}

\aboutMeFrame

\begin{frame}{Disclaimer}
	\begin{itemize}
		\item Diese Folien sind für Tutanden der Veranstaltung \enquote{Grundbegriffe der Informatik} am KIT. Die Weitergabe dieser Folien an Dritte oder eine Veröffentlichung ist nicht gestattet!
		\item Die Folien wurden mit größter Sorgfalt erstellt. Trotzdem ist es wahrscheinlich, dass noch inhaltliche oder formale Fehler enthalten sind. Für die Klausur maßgeblich sind ausschließlich die offiziellen Materialien der Lehrveranstaltung.
	\end{itemize}
\end{frame}

%\begin{headframe}[Das Genie beherrscht das Chaos \\ \forcenewline Oder...]
%	Hashing
%\end{headframe}

\begin{frame}{Grundbegriffe ?}
	\begin{itemize}
		\item Cloud
		\item Künstliche Intelligenz
		\item Big Data
		\item App
		\item Roboter
		\item Informatik
	\end{itemize}

	\pause
	Das alles sind \emph{keine} Grundbegriffe!\\
	Was erwartet euch also?
\end{frame}

\begin{frame} {Grundbegriffe !}
	\begin{itemize}
		\item Menge, Relation, Funktion
		\item Wort, Sprach, Grammatik
		\item Graph
		\item Algorithmus
		\item Turingmaschine
		\item Logik
		\item Prozessor
		\item ...
	\end{itemize}
\end{frame}

\begin{frame}[t]{Tutorium}
	\begin{itemize}[<+->]
		\item Inhalte der Vorlesung erläutern und anwenden
		\item Beispiele, Aufgaben (ihr seid dran!)
		\item Hier sollt ihr Fragen stellen!
		\item \textbf{Kein} Ersatz für die Vorlesung.\\
			  Themen, die nicht im Tutorium dran kommen, sind trotzdem wichtig!
		\item Hier weniger formales Vorgehen als in der Vorlesung. Das heißt nicht, dass das formale vorgehen nicht wichtig ist.
	\end{itemize}
	\YesQuestion{Kriegt ihr hier Antworten?}
	\YesQuestionE{Wird Plagiarismus mit dem Tode bestraft?}{Natürlich! \smiley}
	\NoQuestionE{Dürft ihr mich siezen?}{Ihr werdet sonst erschossen. \smiley}
\end{frame}

\begin{frame}{Übungsblätter}
	Handschriftlich, mit Deckblatt (Tutoriums-Nummer!) und getackert rechtzeitig	abgeben! \\ \pause
	Keine Plagiate, keine Gruppenabgabe! \\[1em] \pause
	Voraussichtlich 6--7 Blätter  Bearbeitungszeit. \\ \pause
	Ausgabe:  Mittwochs, 14 tägig
	Abgabe:   Donnerstag, 16 Uhr, ca \realTilde 2 Wochen später
	Rückgabe: Hier im Tutorium (Nach 2 Wochen Abholung beim Dozenten)
\end{frame}

\begin{frame}{Übungsblätter}
	Noch ein paar spezielle Hinweise:
	\begin{itemize}
		\item Text in Schwarz oder Blau (oder andere dunkle Farbe $\notin \{\text{Rot}, \text{Grün}, \text{Weiß}\}$ )
		\item Links und rechts bitte jeweils einen Rand frei lassen.
		\item Aufgabennummer deutlich kenntlich machen
		\item Für (sehr) unsaubere Abgaben ist Punktabzug möglich!
	\end{itemize}
\end{frame}

\begin{frame} {Leistungen}
	Wir sind Teil der \textbf{Orientierungsprüfung}!
	\begin{itemize}
		\item Übungsschein:\\ 
			50\% aller erreichbarer Punkte hinreichend\\
			Wichtig: Versuchen im ersten Semester (da im Zweiten nicht angeboten)!
		\item Klausur \\
			Teilnahme auch ohne Übungsschein möglich.
			Versuch spätestens im 2. Semester, bestehen im 3.
	\end{itemize}

	\pause
	\DependsQuestionE{Ich studiere \enquote{Europäische Kultur und Ideengeschichte},\\ brauche ich den Übungsschein?}{Es kann nie schaden, die Übungen zu bearbeiten.}
\end{frame}

\begin{frame}
	\frametitle{Kontaktmöglichkeiten}
	
	\begin{block}{Tutorium}
		\begin{itemize}
			\item Mail: \mymail
			\item Folien bekommt ihr im ILIAS.
		\end{itemize} 
	\end{block}
	\pause
	
	\begin{block}{Vorlesung}
		\begin{itemize}
			\item Forum: ILIAS
			\item Bitte alle fachlichen und allgemein organisatorischen Fragen im Forum!
			\item Dozent: sebastian.stueker@kit.edu (Bitte immer Name und Matrikelnummer mit angeben!)
		\end{itemize}
	\end{block}
\end{frame}

\begin{frame}
	\frametitle{Ressourcen}
	
	\begin{itemize}
		\item Vorlesungsfolien
		\item Skript
		\item Archiv (gbi.ira.uka.de)\\
			  Alte Klausuren, Übungsblätter, Lösungen
	\end{itemize}

	\bigskip
	Sehr gute Inhalte, Erklärungen, Übungen findet man auch bei
	\begin{itemize}
		\item EDX (edx.org)
		\item Coursera (coursera.org)
	\end{itemize}
	
\end{frame}

\section{Information}
\begin{frame}
	\frametitle{Nachricht, Information, ...}
	
	\begin{block}{Nachricht}
		      Mitteilung, bei der vom Medium und den Einzelheiten der Signale abstrahiert wird.
	\end{block}
	
	\pause
	\begin{block}{Information}
		      Bedeutung, die einer Nachricht zugeordnet wird (kontextabhängig!).
	\end{block}
	
	\pause
	\begin{block}{Informationsgehalt}
		Bei gleicher Wahrscheinlichkeit: Anzahl der Elemente... \\
		      Naturalis: log$_{e}$ \\
		      Hartley: log$_{10}$ \\
		      Shannon: log$_{2}$ \\[1em]
		Mehr dazu in TGI
	\end{block}

\end{frame}


\section{Mengen}

%TODO: Handout mode
\mycomment{
\begin{frame}
	\frametitle{Problem}
	
	\includegraphics[width=\linewidth]{moviesActors.jpg}
\end{frame}
}

\begin{frame}
	\frametitle{Problem}
	Wir haben \enquote{ein Universum} an Elementen (Filme, Serien, Schauspieler):\\[0.5em]
	
	$U$ enthält Sherlock, Benedict Cumberbatch, Lea Thompson, Martin Freeman, The Imitation Game, Mark Gatiss, Christopher Lloyd, Crispin Glover, Zurück in die Zukunft, Michael Fox, Keira Knightley, \dots  
\end{frame}

\begin{frame}
	\frametitle{Mengen}
	Mengen sind eines der Grundelemente der Mathematik. \\[1em]
	Sammelt bitte 5 Minuten lang \\(alleine / mit eurem Nebensitzer / in Kleingruppen)\\ alles, was ihr über Mengen bereits wisst.
\end{frame}

\begin{frame}
	\frametitle{Mengen}

	\begin{block}{Definition}
		Eine \textbf{Menge} $M$ ist eine Ansammlung verschiedener Objekte. Ein Objekt aus der Menge nennt man ein \textbf{Element} der Menge. Man schreibt
		$$m \in M \qquad M = \{m_1, m_2, m_3 \} \qquad M = \{m \mid m \geq 0\} \qquad M = \emptyset = \{\} $$
		Die Reihenfolge der Aufzählung ist dabei irrelevant, Elemente kommen nicht doppelt vor. Die leere Menge $\emptyset$ enthält keine Elemente.
	\end{block}
	\pause
	
	\begin{block}{Beispiel}
		Wichtige Mengen sind
		$$\N, \Z, \Q, \R, \C\qquad \N_+, \N_0$$
		Es gilt: $$-5 \in \Z\qquad -5 \notin \N\qquad \{2, 1, 3, 1, 4\} = \{1, 2, 3, 4\}$$
	\end{block}
	
\end{frame}

\begin{frame}
	\frametitle{Teilmengen}
	
	\begin{block}{Definition}
		Die Anzahl der Elemente in einer endlichen Menge (Kardinalität) bezeichnet man mit $\setsize{M}$. Es gilt $$ \setsize{M \cup N} \ = \ \setsize{M} + \setsize{N} - \setsize{M \cap N} $$
	\end{block} 
	\pause

	\begin{block}{Definition}
		Eine Menge $N$ ist eine \textbf{Teilmenge} von $M$, wenn jedes Element aus N auch in M enthalten ist.
		$$ N \subseteq M \iff \forall n \in N : n \in M$$
	\end{block}
	\pause
	
	\begin{block}{Definition}
		Zwei Mengen $N$ und $M$ sind \textbf{gleich}, wenn sie die gleichen Elemente enthalten.
		$$ N = M \iff N \subseteq M \ \text{und} \ M \subseteq N$$
	\end{block}
	
\end{frame}

\begin{frame}
	\frametitle{Teilmengen}
	
	\begin{block}{Beispiel}
		
	\begin{align*}
		\setsize{\{1,2,3,2,1\}} &= \only<2->{3} \\
		\setsize{ \emptyset } &= \only<3->{0} \\
		\{1,2\} \only<1-3|handout:0>{\ &\mathrel{?} \ } \only<4->{&\subseteq}  \{1,2,3\} \\
		\{1,2\} \only<1-4|handout:0>{\ &\mathrel{?} \ } \only<5->{&\nsubseteq} \{\text{Hund}, \text{Katze}, \text{Maus}\}
	\end{align*}
	\end{block} 

	\only<6->{
		\begin{block}{Lemma}
			Es gilt: $$ N \subseteq M \iff N \setminus M = \emptyset $$
			
		\end{block} 

	}
\end{frame}

\begin{frame}{Mengengleichheit: Beispiel}
	\begin{itemize}
		\item<1-> Sei $ A $ und $M$ beliebige Mengen. Zeigen Sie, dass gilt 
		\begin{align*}
		A &=  \underbrace{ \left(A \setminus M \right)}_{T_1} \cup  \underbrace{ \left(A \cap M \right)}_{T_2}  
		\end{align*}		 
		\item<2-5>
		Richtung: $ A \subseteq T_1 \cup T_2 $ 
		\only<3-5>{\item<3-5> Wähle $x\in A$ und wende Fallunterscheidung an
			\item<4-5> Fall 1: Ist $x\in M$ so gilt $x\in A $ und $x\in M$, und damit $x\in A\cap M  = T_2 $
			\item<5-5> Fall 2 : Ist $x\notin M$ so gilt $x\in T_1$ da $T_1 = \left\{ x \in A \text{ und } x\notin M \right\} $  }
		\only<6->{\item<6-> Richtung : $T_1 \cup T_2 \subseteq A $ 
			\item<7-> Wähle $x\in T_1 \cup T_2$. Dies bedeutet $x\in T_1 $ oder $x\in T_2$. 
			\item<8-> Fall 1: $x\in T_1$. Aus Definition folgt $x\in A$
			\item<9-> Fall 2: $x\in T_2$. Somit $x\in A$ und $x\in M$.}
	\end{itemize}
\end{frame}

\begin{frame}
	\frametitle{Zurück zu unserem Problem}
	$U = \{ $ Sherlock, Benedict Cumberbatch, Lea Thompson, Martin Freeman, The Imitation Game, Mark Gatiss, Christopher Lloyd, Crispin Glover, Zurück in die Zukunft, Michael Fox, Keira Knightley, \dots $ \} $  \\[0.5em]
	
	\pause
	Ordnen wir diese in eine Teilmenge $M$ für Filme/Serien und jeweils eine Teilmenge $A_m$ für die Schauspieler eines Films $m \in M$. \\[0.5em]
	
	\pause
	$M = \{$ The Imitation Game, Sherlock, Zurück in die Zunkunft $\}$ \\[0.3em]
	$A_{Sherlock} = \{$ Benedict Cumberbatch, Martin Freeman, Mark Gatiss $\}$ \\
	$A_{Imitation Game} = \{$ Benedict Cumberbatch, Keira Knightley $\}$ \\
	$A_{BTTF} = \{$ Michael J. Fox, Christopher Lloyd, Lea Thompson, Crispin Glover $\}$
	
\end{frame}

\begin{frame}
	\frametitle{Schnitt und Vereinigung}
	
	\begin{block}{Definition}
		Sind $M$ und $N$ zwei Mengen, so definiert man
		$$M \cap N = \{x \mid x \in M \text{ und } x \in N\} \qquad M \cup N = \{x \mid x \in M \text{ oder } x \in N\} $$
		als den \textbf{Durchschnitt} und die \textbf{Vereinigung}.\\[1em] 
		\pause
		Zwei Mengen $M$, $N$ heißen \textbf{disjunkt}, wenn ihr Durchschnitt leer ist, sie also keine gemeinsamen Elemente besitzen. $$M \cap N = \emptyset$$
	\end{block}

	\pause
	\begin{block}{Beispiel}
		$$ \{1,2\} \cup \{2,3\} = \{1,2,3\} \qquad \{1,2\} \cap \{2,3\} = \{2\} $$
	\end{block}

\end{frame}

\begin{frame}
	\frametitle{Schnitt und Vereinigung}
	\begin{block}{Definition}
		Seien $A$ und $B$ zwei beliebige Mengen, so gilt $$ A\setminus B = \left\{ x\in A \text{ und } x\notin  B  \right\} $$ 
	\end{block}
	
	\pause

	\begin{block}{Weitere Beispiele}
		\begin{align*}
			A \cup \emptyset \only<3->{ &= A }  \\
			A \cap \emptyset \only<4->{ &= \emptyset }\\
			\N_+ \cup \{0\} \only<5->{ &= \N_0} \\
		\end{align*}
	\end{block}

\end{frame}

\begin{frame}
	\frametitle{Eine Menge Mengen...}
	\begin{block}{Aufgabe}
		Es seien $A = \{1, 2\}, B = \{3\}, C = \{1, 3\}  \subseteq M = \{1, 2, 3\}$ Mengen.\\
		Man berechne folgende Mengen:
		\begin{align*}
		A \cup B &= \only<2->{ \{1, 2, 3\} }  \\
		A \cap C &= \only<3->{ \{1\} }\\
		A \setminus C &= \only<4->{ \{2\} }\\
		B \setminus A &= \only<5->{ \{3\} }\\
		A \cup (B \setminus C) &= \only<6->{ \{1, 2\} }\\
		C &= \{1, 3\} \\
		(A \setminus C) \cup B &= \only<7->{ \{2, 3\} }\\
		A \cap B &= \only<8->{ \emptyset }
		\end{align*}
	\end{block}
\end{frame}

\section{Potenzmengen}
\begin{frame}
	\frametitle{... in einer Menge! (Potenzmengen)}
	\begin{block}{Definition}
		Die \textbf{Potenzmenge} $2^M$ oder auch $\Pot (M)$ ist die Menge aller möglicher Teilmengen von $M$. Es gilt also 
		\begin{align*}
			2^M = \{N \mid N \subseteq M\}
		\end{align*}
	\end{block}
	\pause
	
	\begin{block}{Beispiel}
		Betrachten wir nun   $M = \left\{ 1,2,0 \right\} $. \\ \pause
		
		Dann gilt 
		\begin{align*}
		2^M &= \left\{ \emptyset, \left\{ 0 \right\}, \left\{ 1 \right\}, \left\{ 2 \right\}, \left\{ 0,1 \right\} , \left\{ 0,2 \right\}, \left\{ 1,2 \right\}, \left\{ 0,1,2 \right\} \right\}
		\end{align*}
		
		Beachte: Es gilt immer $M \in 2^M \ \text{und} \ \emptyset \in 2^M$
	\end{block}
	
\end{frame}

\begin{frame}
	\frametitle{Potenzmengen}
	\begin{block}{Aufgabe}
		Wie viele Elemente enthält $2^M$? \\[0.5em]
		
		\pause
		$2^{\mid M \mid }$
	\end{block}
	
	\pause
	\begin{block}{Aufgabe}
		Geben Sie eine Abbildung $\phi \colon 2^{M} \longrightarrow 2^{M}$ so an,
		dass für jedes $L \in 2^{M}$ und für jedes $w \in M$ gilt:
		\begin{equation*}
			w \in L \text{ genau dann, wenn } w \notin \phi(L).
		\end{equation*}
		
		\pause
		\begin{align*}
			\phi \colon 2^{M} &\longrightarrow 2^{M},\\
			L &\mapsto M \setminus L.
		\end{align*}
	\end{block}

\end{frame}

\begin{frame}
	\frametitle{Zurück zu unserem Problem}
	Ordnen wir diese in eine Teilmenge $M$ für Filme/Serien und jeweils eine Teilmenge $A_m$ für die Schauspieler eines Films $m \in M$. \\[0.5em]
	$M = \{$ The Imitation Game, Sherlock, Zurück in die Zunkunft $\}$ \\[0.3em]
	$A_{Sherlock} = \{$ Benedict Cumberbatch, Martin Freeman, Mark Gatiss $\}$ \\
	$A_{Imitation Game} = \{$ Benedict Cumberbatch, Keira Knightley $\}$ \\
	$A_{BTTF} = \{$ Michael J. Fox, Christopher Lloyd, Lea Thompson, Crispin Glover $\}$ \\[2em]
	
	\pause
	$A_{Sherlock} \cap A_{Imitation Game} = \{ \text{Benedict Cumberbatch} \}$
	
\end{frame}

\section{Paare}
\begin{frame}
	\frametitle{Paare}
	\begin{block}{Definition}
		Seien $A$ und $B$ zwei Mengen und $a \in A$, $b \in B$.\\
		$$(a, b)$$ heißt \textbf{Paar} mit der ersten Komponente $a$ und der zweiten Komponente $b$.\\[1em]
		\pause
		In Paaren können Elemente mehrfach vorkommen, und die Reihenfolge der Elemente ist wichtig.\\
	\end{block}

	\pause
	\begin{block}{Beispiel}
		$$ (KI, T) = (KI, T) \qquad (KI, T) \neq (T, KI) \qquad (1, 1) $$
	\end{block}
\end{frame}

\begin{frame}
	\frametitle{Paare}
	Das Konzept der Paare lässt sich auf das Konzept der Mengen zurückführen.
	
	\begin{block}{Aufgabe}
		Gegeben sei die Menge $M = \{m_1, m_2\}$.\\
		Wie kann man die Paare $(m_1, m_2)$ und $(m_2, m_1)$ eindeutig darstellen, nur unter Verwendung von Mengen und $m_1, m_2$?
	\end{block}
	\pause
	\begin{block}{Lösung}
		Wir definieren: \\
		$(m_1, m_2) := \{m_1, m_2, \{m_1\}\}$ und $(m_2, m_1) := \{m_1, m_2, \{m_2\}\}$
	\end{block}
\end{frame}


\begin{frame}
	\frametitle{}
	
	\begin{block}{Was ihr nun wissen solltet}
		\begin{itemize}
			\item wie man mit Mengen umgeht
			\item wie man mit noch mehr Mengen umgeht
			\item was Paare sind
		\end{itemize}
	\end{block}

	\begin{block}{Was nächstes Mal kommt}
		\begin{itemize}
			\item wie man mehr als zwei Elemente geordnet zusammenbringt (Tupel)
			\item wie man mit Relationen Ordnung in das Chaos bringt
			\item [...] und noch vieles mehr!
			%TODO: Fg color := black for \item []-titles!
		\end{itemize}
	\end{block}
\end{frame}

%% Letzte Seite 
\lastframe{0.35}{30}{xkcd/good_code.png}{http://xkcd.com/844}

\slideThanks

\end{document}
